\title{QLoRA: Efficient Finetuning of Quantized LLMs}

\begin{abstract}
We introduce \textsc{QLoRA}, a memory-efficient finetuning method that enables the finetuning of a 65B parameter model on a single 48GB GPU while maintaining the performance of full 16-bit finetuning tasks. \textsc{QLoRA} backpropagates gradients through a frozen, 4-bit quantized pretrained language model into Low Rank Adapters~(LoRA). Our top-performing model series, named \textbf{Guanaco}, surpasses all previously publicly released models on the Vicuna benchmark, achieving 99.3\% of ChatGPT’s performance with just 24 hours of finetuning on a single GPU. \textsc{QLoRA} introduces several innovations to reduce memory usage without compromising performance: (a) 4-bit NormalFloat (NF4), a novel data type that is theoretically optimal for normally distributed weights, (b) Double Quantization, which minimizes average memory usage by quantizing the quantization constants themselves, and (c) Paged Optimizers to handle memory spikes efficiently. We apply \textsc{QLoRA} to finetune over 1,000 models, conducting an extensive analysis of instruction-following and chatbot capabilities across 8 instruction datasets, various model architectures (LLaMA, T5), and model sizes that would be impractical to finetune conventionally (e.g., 33B and 65B parameter models). Our findings indicate that finetuning with QLoRA on a small, high-quality dataset achieves state-of-the-art outcomes even with smaller models than previous SoTA. We present a comprehensive evaluation of chatbot performance based on both human and GPT-4 assessments, demonstrating that GPT-4 evaluations offer a cost-effective and reliable alternative to human evaluation. Moreover, we reveal that existing chatbot benchmarks do not reliably measure chatbot performance levels. A targeted analysis highlights the specific shortcomings of \textbf{Guanaco} relative to ChatGPT. We publicly release all models and code, including CUDA kernels for 4-bit training.\footnote{\url{https://github.com/artidoro/qlora} and \url{https://github.com/TimDettmers/bitsandbytes}}
\end{abstract}

\section{Introduction}

Finetuning large language models (LLMs) is an extremely effective approach to enhance their capabilities \citep{min2021metaicl, wei2021finetuned, ouyang2022training, wang2022super, wang2022self, liu2022few}, as well as to incorporate desired behaviors or eliminate unwanted ones \citep{ouyang2022training, askell2021general, bai2022training}. Nonetheless, finetuning very large models remains prohibitively costly; for instance, standard 16-bit finetuning of a LLaMA 65B parameter model~\citep{touvron2023llama} demands over 780 GB of GPU memory. Although recent quantization methods can shrink the memory footprint of LLMs \citep{dettmers2022llm, dettmers2022case, frantar2022gptq, xiao2022smoothquant}, these techniques are only suitable for inference and fail during training \citep{wortsman2023stable}. 

We present, for the first time, the feasibility of finetuning a quantized 4-bit model without any loss in performance. Our approach, \textsc{QLoRA}, applies an innovative high-precision quantization method to compress a pretrained model to 4-bit precision, then introduces a small set of learnable Low-rank Adapter weights \citep{hu2021lora} that are optimized by backpropagating gradients through the quantized weights. 

\textsc{QLoRA} cuts down the average memory usage for finetuning a 65B parameter model from over 780GB of GPU memory to under 48GB, without compromising runtime efficiency or predictive accuracy when compared to a fully 16-bit finetuned baseline. This represents a major advancement in the accessibility of LLM finetuning, enabling the largest publicly available models to be finetuned on a single GPU. Utilizing \textsc{QLoRA}, we develop the \textbf{Guanaco} family of models, where the second-best model achieves 97.8\% of ChatGPT’s performance on the Vicuna~\citep{vicuna2023} benchmark, while requiring less than 12 hours of training on a single consumer GPU; with a single professional GPU and 24 hours of training, our largest model reaches 99.3\%, effectively closing the performance gap to ChatGPT on the Vicuna benchmark. When deployed, our smallest \textbf{Guanaco} model (7B parameters) uses only 5 GB of memory and outperforms a 26 GB Alpaca model by over 20 percentage points on the Vicuna benchmark (Table~\ref{tbl:vicuna_eval}).

\textsc{QLoRA} presents several innovations aimed at minimizing memory consumption without compromising performance: (1) \textbf{4-bit NormalFloat}, an information-theoretically optimal quantization format for data with a normal distribution, which delivers superior empirical results compared to 4-bit Integers and 4-bit Floats. (2) \textbf{Double Quantization}, a technique that further quantizes the quantization constants themselves, resulting in an average savings of approximately 0.37 bits per parameter (around 3 GB for a 65B parameter model). (3) \textbf{Paged Optimizers}, which leverage NVIDIA unified memory to prevent memory spikes during gradient checkpointing that happen when processing mini-batches with long sequence lengths. These innovations are integrated into a finely tuned LoRA approach incorporating adapters at every network layer, thus almost entirely eliminating the accuracy compromises seen in previous work.

The efficiency of \textsc{QLoRA} allows for a comprehensive investigation into instruction finetuning and chatbot performance across model scales that would be unfeasible with standard finetuning due to memory limitations. Consequently, we train over 1,000 models spanning various instruction tuning datasets, model architectures, and sizes ranging from 80M to 65B parameters. Beyond demonstrating that \textsc{QLoRA} matches 16-bit precision performance (\S\ref{sec:comp_to_std}) and enabling the training of a cutting-edge chatbot, \textbf{Guanaco}, (\S\ref{sec:pushing}), we also explore patterns in the trained models. Firstly, we observe that data quality is significantly more critical than dataset size; for example, a 9k sample dataset (OASST1) outperformed a 450k sample dataset (FLAN v2, subsampled) in chatbot performance, despite both targeting instruction-following generalization. Secondly, we reveal that achieving strong results on the Massive Multitask Language Understanding (MMLU) benchmark does not guarantee strong performance on the Vicuna chatbot benchmark and vice versa—highlighting that the appropriateness of the dataset is more influential than its size for a given task.

Additionally, we conduct a comprehensive evaluation of chatbot performance utilizing assessments from both human evaluators and GPT-4. Our approach employs a tournament-style benchmarking framework where different models face off in matches to generate the optimal response for a specific prompt. The victor of each match is determined by either GPT-4 or human judges. The outcomes of these tournaments are compiled into Elo scores~\citep{elo1967proposed, elo1978rating}, which are then used to rank the chatbot performances. Our findings indicate that GPT-4 and human evaluations largely concur on the ranking of model capabilities in these tournaments, although there are certain cases of notable disagreement. Therefore, we emphasize that although model-based evaluation offers a cost-effective alternative to human annotation, it also entails some degree of uncertainty.

We complement our benchmark results on chatbots with a qualitative examination of the \textbf{Guanaco} models. This qualitative analysis sheds light on success and failure instances that are not reflected in the quantitative metrics.

To support further research, we release all model-generated responses along with annotations from both human evaluators and GPT-4. We also open-source our codebase and CUDA kernels, integrating our techniques into the Hugging Face transformers library \citep{wolf2019huggingface} to ensure broad accessibility. Furthermore, we provide a set of adapters for models of sizes 7B, 13B, 33B, and 65B, each trained on 8 distinct instruction-following datasets, resulting in a total of 32 publicly available fine-tuned models.

\section{Background}
\paragraph{Block-wise k-bit Quantization} Quantization refers to the process of converting data from a representation containing more information to one with reduced information content. This usually involves transforming data types with higher bit-widths to those with fewer bits, such as converting 32-bit floating-point numbers to 8-bit integers. To fully utilize the representational range of the lower-bit data type, the original data is typically normalized by dividing by the absolute maximum of the input tensor's elements. For instance, when quantizing a 32-bit floating-point tensor into an 8-bit integer tensor within the range $[-127, 127]$, we have:

\begin{equation}
    \mathbf{X}^{ \text{Int8}} = \text{round}\left( \frac{127}{{\text{absmax}}(\mathbf{X}^{{\text{FP32}}})}\mathbf{X}^{{\text{FP32}}} \right) = \text{round}(c^{\text{FP32}}\cdot \mathbf{X}^{{\text{FP32}}}),
\end{equation}

where $c$ denotes the {\em quantization constant} or {\em quantization scale}. The reverse process, dequantization, is given by:

\begin{equation}
    \text{dequant}(c^{\text{FP32}}, \mathbf{X}^{\text{Int8}}) = \frac{\mathbf{X}^{\text{Int8}}}{c^{\text{FP32}}} = \mathbf{X}^{\text{FP32}}
\end{equation}

The issue with this method is that when an input tensor contains a value with a large magnitude (i.e., an outlier), the quantization bins—specific bit patterns—are inefficiently used, resulting in some bins having few or no quantized values. To address the outlier problem, a common strategy is to divide the input tensor into smaller blocks, each quantized independently with its own quantization constant $c$. Formally, the input tensor $\mathbf{X}\in \mathbb{R}^{b\times h}$ is split into $n$ consecutive blocks of size $B$ by flattening $\mathbf{X}$ and segmenting the linear array into ${n = ({b\times h} )/{B} }$ segments. Each block is quantized separately using Equation~1, yielding a quantized tensor along with $n$ quantization constants $c_i$.

\paragraph{Low-rank Adapters} Low-rank Adapter (LoRA) finetuning \citep{hu2021lora} is a technique aimed at reducing memory consumption by training a small set of additional parameters called adapters, while keeping the original model parameters fixed. During stochastic gradient descent, gradients flow through the fixed pretrained weights to the adapters, which are then updated to minimize the loss function. LoRA enhances a linear transformation by adding an extra factorized projection. For a projection ${\mathbf{X} \mathbf{W} = \mathbf{Y}}$ where $\mathbf{X}\in  \mathbb{R}^{b\times h}$ and $\mathbf{W}\in  \mathbb{R}^{h\times o}$, LoRA computes:

\begin{equation}
    \mathbf{Y} = \mathbf{X}\mathbf{W} + s\mathbf{X}\mathbf{L}_1\mathbf{L}_2,
\end{equation}

in which $\mathbf{L}_1\in  \mathbb{R}^{h\times r}$, $\mathbf{L}_2\in  \mathbb{R}^{r\times o}$, and $s$ is a scalar.

\paragraph{Memory Requirement of Parameter-Efficient Finetuning} A key aspect to consider is the memory usage of LoRA during training, specifically regarding the quantity and size of adapters deployed. Given LoRA's extremely small memory footprint, it is feasible to utilize a greater number of adapters to enhance performance without substantially increasing the total memory consumption. Although LoRA was developed as a Parameter Efficient Finetuning (PEFT) approach, the majority of memory used in LLM finetuning is attributed to activation gradients rather than the LoRA parameters themselves. For instance, when training a 7B LLaMA model on FLAN v2 with a batch size of 1, and LoRA weights corresponding to approximately 0.2\% of the original model parameters\citep{hu2021lora,liu2022few}, the memory occupied by LoRA input gradients is 567 MB, whereas the LoRA parameters only require 26 MB. Employing gradient checkpointing~\citep{chen2016training} reduces input gradients to about 18 MB per sequence on average, making them more memory demanding than all LoRA parameters combined. By contrast, the 4-bit base model utilizes 5,048 MB of memory. This underscores the significance of gradient checkpointing and indicates that aggressively minimizing LoRA parameter size offers only limited memory savings. Consequently, it is possible to incorporate more adapters without notably increasing the overall training memory requirements (refer to Appendix~\ref{app:memory} for a comprehensive analysis). As elaborated later, this capability is vital for achieving full 16-bit precision performance.

\section{\textsc{QLoRA} Finetuning}

\textsc{QLoRA} attains high-accuracy 4-bit finetuning through two key methods we introduce—4-bit NormalFloat (NF4) quantization and Double Quantization. Furthermore, we present Paged Optimizers, designed to prevent memory spikes during gradient checkpointing from triggering out-of-memory errors, which have historically made finetuning large models on a single machine challenging.

\textsc{QLoRA} employs a single low-precision storage format, typically 4-bit in our implementation, alongside a computation data type, usually BFloat16. In practice, this means whenever a \textsc{QLoRA} weight tensor is accessed, we first dequantize it to BFloat16 before performing the matrix multiplication using 16-bit precision.

We will now outline the elements of \textsc{QLoRA} and then provide a formal definition of the method.

\paragraph{4-bit NormalFloat Quantization} The NormalFloat (NF) data format is built upon Quantile Quantization \citep{dettmers2022optimizers}, an information-theoretically optimal scheme that ensures each quantization bin contains an equal number of values from the input tensor. Quantile quantization operates by estimating the quantiles of the input tensor via the empirical cumulative distribution function.

A key drawback of quantile quantization is the computational cost associated with estimating quantiles. Therefore, fast approximation methods, such as SRAM quantiles~\citep{dettmers2022optimizers}, are utilized to approximate them. However, due to the approximate nature of these methods, the data type suffers from significant quantization errors for outliers, which are often the most critical values.

Such costly quantile estimation and approximation inaccuracies can be avoided when input tensors originate from a distribution fixed up to a quantization constant. In these scenarios, input tensors share identical quantiles, making exact quantile estimation computationally viable.

Pretrained neural network weights typically follow a zero-centered normal distribution with a standard deviation $\sigma$ (refer to Appendix~\ref{app:normality}). By scaling $\sigma$, we can adjust all weights to conform to a single fixed distribution that precisely fits within the range of our chosen data type. For our purposes, we define this arbitrary range as $[-1, 1]$. Consequently, both the quantiles of the data type and those of the neural network weights must be normalized to this interval.

To derive the information-theoretically optimal data type for zero-mean normal distributions with any standard deviation $\sigma$ constrained to the $[-1, 1]$ range, we proceed as follows: (1) calculate the $2^k+1$ quantiles of the theoretical $N(0, 1)$ distribution to generate a $k$-bit quantile quantization data type tailored for normal distributions, (2) normalize the values of this data type into the $[-1, 1]$ range, and (3) quantize an input weight tensor by normalizing it into $[-1, 1]$ through absolute maximum rescaling.

After aligning the weight tensor range with that of the data type, standard quantization can be applied. The third step is equivalent to adjusting the weight tensor's standard deviation to equal that of the $k$-bit data type. More precisely, we define the $2^k$ data type values $q_i$ as:

\begin{equation}
q_i  = \frac{1}{2}\left( Q_X\left(\frac{i}{2^k+1}\right) + Q_X\left(\frac{i+1}{2^k+1}\right)\right),
\end{equation}

where $Q_X(\cdot)$ denotes the quantile function of the standard normal distribution $N(0,1)$. A challenge with symmetric k-bit quantization is that it lacks an exact zero representation, which is crucial for quantizing padding and other zero-valued elements without error. To guarantee a discrete zero point at $0$ and to fully utilize all $2^k$ levels for a k-bit datatype, we construct an asymmetric datatype by estimating quantiles $q_i$ over two ranges: $2^{k-1}$ for the negative side and $2^{k-1}+1$ for the positive side. We then merge these sets of $q_i$ and eliminate one of the two zeros that appear in both sets. We call the resulting datatype, which has an equal expected count of values per quantization bin, \emph{k-bit NormalFloat} (NFk), as it is information-theoretically optimal for zero-centered normally distributed data. Precise values for this datatype are provided in Appendix~\ref{app:nf4}.

\paragraph{Double Quantization{}} 
We propose \emph{Double Quantization} (DQ), a method of quantizing the quantization constants themselves to achieve further memory savings. Although small block sizes are necessary for accurate 4-bit quantization~\citep{dettmers2022case}, they cause significant memory overhead. For instance, using 32-bit constants with a block size of 64 for $\mathbf{W}$ results in an average addition of $32/64 = 0.5$ bits per parameter due to quantization constants. Double Quantization aims to reduce this memory overhead.

In detail, Double Quantization treats the quantization constants $c_2^{\text{FP32}}$ from the first quantization as inputs to a second quantization step. This yields quantized constants $c_2^{\text{FP8}}$ and a new set of quantization constants $c_1^{\text{FP32}}$ for this second level. We employ 8-bit floating-point numbers with a block size of 256 for this second quantization, as no accuracy loss is observed at 8 bits, consistent with findings by \citet{dettmers2022case}. Since the $c_2^{\text{FP32}}$ values are positive, we first subtract their mean to center them around zero, allowing for symmetric quantization. For a block size of 64, this approach reduces the memory cost per parameter from $32/64 = 0.5$ bits to $8/64 + 32/(64 \cdot 256) = 0.127$ bits on average, yielding a saving of 0.373 bits per parameter.

\paragraph{Paged Optimizers} leverage the NVIDIA unified memory~\footnote{\tiny \url{https://docs.nvidia.com/cuda/cuda-c-programming-guide}} capability, which performs automatic page-by-page data transfers between the CPU and GPU to ensure error-free GPU computations in situations where the GPU occasionally experiences out-of-memory conditions. This mechanism operates similarly to conventional memory paging between CPU RAM and disk storage. We utilize this feature to allocate paged memory for the optimizer states, which are then automatically moved to CPU RAM when GPU memory is insufficient and brought back to GPU memory when required during the optimizer update phase.

\paragraph{\textsc{QLoRA}.} Building on the components outlined earlier, we define \textsc{QLoRA} for a single linear layer within the quantized base model incorporating a single LoRA adapter as follows:

\begin{equation}
    \mathbf{Y}^{\text{BF16}} =  \mathbf{X}^{\text{BF16}} \text{doubleDequant}(c_1^{\text{FP32}}, c_2^{\text{k-bit}}, \mathbf{W}^{\text{NF4}}) + \mathbf{X}^{\text{BF16}}\mathbf{L}^{\text{BF16}}_1\mathbf{L}^{\text{BF16}}_2,
\end{equation}

where doubleDequant$(\cdot)$ is defined by:

\begin{equation}
    \text{doubleDequant}(c_1^{\text{FP32}}, c_2^{\text{k-bit}}, \mathbf{W}^{\text{k-bit}}) = \text{dequant}(\text{dequant}(c_1^{\text{FP32}}, c_2^{\text{k-bit}}), \mathbf{W}^{\text{4bit}}) = \mathbf{W}^{\text{BF16}},
\end{equation}

Here, we adopt NF4 for $\mathbf{W}$ and FP8 for $c_2$.

A blocksize of 64 is selected for $\mathbf{W}$ to enhance quantization accuracy, while a blocksize of 256 is chosen for $c_2$ to optimize memory usage.

During parameter updates, only the gradients with respect to the adapter weights $\frac{\partial E}{\partial \mathbf{L}_i}$ are required, not the gradients for the 4-bit weights $\frac{\partial E}{\partial \mathbf{W}}$. Nevertheless, computing $\frac{\partial E}{\partial \mathbf{L}_i}$ requires evaluating $\frac{\partial \mathbf{X}}{\partial \mathbf{W}}$, which proceeds according to equation~(5), involving dequantization from stored format $\mathbf{W}^{\text{NF4}}$ to the computation data type $\mathbf{W}^{\text{BF16}}$ in order to calculate the derivative $\frac{\partial \mathbf{X}}{\partial \mathbf{W}}$ with BFloat16 precision.

In summary, \textsc{QLoRA} utilizes a single storage data type (commonly 4-bit NormalFloat) alongside a computation data type (16-bit BrainFloat). To carry out the forward and backward passes, the storage data type is dequantized into the computation data type, but weight gradients are computed solely for the LoRA parameters, which are represented using 16-bit BrainFloat.

\section{QLoRA Compared to Standard Finetuning}
\label{sec:comp_to_std}

We have explained the mechanism of QLoRA and its capability to drastically lower the memory needed for model finetuning. The central question now is whether QLoRA can achieve performance comparable to full-model finetuning. Additionally, we aim to evaluate the various elements of QLoRA, such as the effect of NormalFloat4 compared to conventional Float4. The subsequent sections detail experiments designed to address these questions.

\paragraph{Experimental setup.} Our study covers three model architectures—encoder, encoder-decoder, and decoder-only—and compares QLoRA against 16-bit adapter finetuning and full finetuning for models up to 3B parameters. Our assessments encompass GLUE \citep{wang2018glue} with RoBERTa-large \citep{liu2019roberta}, Super-NaturalInstructions (TKInstruct) \citep{wang2022super} using T5 \citep{t5}, and 5-shot MMLU \citep{hendrycksmeasuring} following finetuning of LLaMA on Flan v2 \citep{longpre2023flan} and Alpaca \citep{alpaca}. To further investigate the benefits of NF4 over other 4-bit data formats, we employ the experimental setup from \citet{dettmers2022case} and evaluate post-quantization zero-shot accuracy and perplexity across various models (OPT~\citep{zhang2022opt}, LLaMA~\citep{touvron2023llama}, BLOOM~\citep{scao2022bloom}, Pythia~\citep{biderman2023pythia}) spanning sizes from 125M to 13B parameters.

Additional specifics for each experimental setup are provided in the results section to enhance clarity. Complete details can be found in Appendix~\ref{app:exp-setup}.

Although paged optimizers are essential to perform 33B/65B \textsc{QLoRA} finetuning on a single 24/48GB GPU, we do not present definitive benchmarks for them because paging occurs primarily during processing of mini-batches with very long sequences, which is infrequent. Nonetheless, we conduct a runtime analysis of paged optimizers on 65B models running on 48GB GPUs, showing that with a batch size of 16, paged optimizers achieve training speeds comparable to standard optimizers. Future research should focus on identifying conditions under which paging-induced slowdowns arise.

\paragraph{Standard LoRA hyperparameters do not replicate 16-bit finetuning performance}

Using the conventional approach of applying LoRA to the query and value attention projection matrices \citep{hu2021lora}, we fail to reproduce full finetuning results for large base models. As illustrated in Figure~\ref{fig:hyper} for LLaMA 7B finetuning on Alpaca, we observe that the most influential LoRA hyperparameter is the total number of LoRA adapters employed. Applying LoRA to all linear layers within transformer blocks is necessary to reach full finetuning performance. Other LoRA hyperparameters, such as the projection dimension $r$, show no significant impact on results (refer to Appendix \ref{app:exp-setup}).

Similarly, we observe that the default hyperparameters for fully finetuned baselines are not optimally tuned. We conduct a hyperparameter search over learning rates ranging from 1e-6 to 5e-5 and batch sizes from 8 to 128 in order to establish robust baselines. The results for finetuning the 7B LLaMA model on Alpaca are presented in Figure~\ref{fig:hyper}.

\paragraph{4-bit NormalFloat outperforms 4-bit Floating Point}

Although the 4-bit NormalFloat (NF4) data type is theoretically optimal from an information perspective, it remains to be verified whether this theoretical advantage translates into practical improvements. We replicate the experimental framework from \citet{dettmers2022case}, evaluating quantized large language models (OPT \citep{zhang2022opt}, BLOOM \cite{scao2022bloom}, Pythia \cite{biderman2023pythia}, LLaMA) of varying sizes (125M to 65B parameters) with different data types on language modeling and a suite of zero-shot tasks. As shown in Figure~\ref{fig:nf4} and Table~\ref{tbl:nf4_ppl}, NF4 yields notably better performance compared to FP4 and Int4, while double quantization effectively reduces memory usage without compromising performance.

\paragraph{k-bit \textsc{QLoRA} achieves parity with 16-bit full finetuning and 16-bit LoRA}

Recent studies have demonstrated that 4-bit quantization for \emph{inference} is achievable, but typically results in performance drops relative to 16-bit precision \citep{dettmers2022case,frantar2022gptq}. This poses the important question of whether such performance loss can be recovered by applying 4-bit adapter finetuning. We investigate this in two different configurations.

The first setup compares full 16-bit finetuning with 4-bit adapter finetuning on RoBERTA and T5 models ranging from 125M to 3B parameters, evaluated on GLUE and the Super-NaturalInstructions datasets. Results displayed in Table~\ref{tab:nf4vsbf16} indicate that adapter methods using 16-bit, 8-bit, and 4-bit precision all reproduce the performance levels of the fully finetuned 16-bit baseline. This demonstrates that the performance degradation caused by quantization can be entirely recovered through adapter finetuning performed after quantization.

For our second experimental configuration, given that fully finetuning models with 11 billion parameters or more demands multiple servers equipped with high-memory GPUs, we proceed to investigate whether 4-bit \textsc{QLoRA} can achieve comparable results to 16-bit LoRA across the 7B to 65B parameter range. To accomplish this, we finetune LLaMA models ranging from 7B to 65B parameters on two instruction-following datasets, Alpaca and FLAN v2, and assess performance on the MMLU benchmark using 5-shot accuracy. The outcomes, presented in Table~\ref{tbl:llama-datatype}, demonstrate that NF4 with double quantization completely recovers the MMLU performance of 16-bit LoRA. Furthermore, we observe that \textsc{QLoRA} employing FP4 falls behind the 16-bit brain float LoRA baseline by roughly one percentage point. This supports our conclusions that (1) \textsc{QLoRA} using NF4 matches the performance of both 16-bit full finetuning and 16-bit LoRA finetuning, and (2) NF4 surpasses FP4 in terms of quantization accuracy.

\paragraph{Summary} Our findings consistently indicate that 4-bit \textsc{QLoRA} with the NF4 data type attains performance equivalent to 16-bit full finetuning and 16-bit LoRA finetuning on academic benchmarks featuring well-established evaluation protocols. We have also demonstrated that NF4 outperforms FP4 and that implementing double quantization does not impair performance. Taken together, this offers strong evidence that 4-bit \textsc{QLoRA} tuning reliably delivers results on par with 16-bit approaches.

Consistent with prior research on quantization \citep{dettmers2022case}, our MMLU and Elo evaluations suggest that, given fixed resources for finetuning and inference, it is advantageous to increase the base model's parameter count while reducing their numerical precision. This underscores the significance of the efficiency advantages provided by \textsc{QLoRA}. Since we detected no performance decline relative to full finetuning in our 4-bit finetuning experiments, it raises an important question about the precise nature of the performance-precision trade-off in QLoRA tuning, which we propose to investigate in future work.

We move forward to explore instruction tuning at scales that would be unfeasible to examine using full 16-bit finetuning on standard academic research hardware.

\section{Advancing the Chatbot State-of-the-art with QLoRA}
\label{sec:pushing}
Having demonstrated that 4-bit \textsc{QLoRA} achieves parity with 16-bit performance across various model sizes, tasks, and datasets, we conduct a comprehensive investigation of instruction finetuning on the largest open-source language models currently accessible for research. To evaluate the effectiveness of instruction finetuning on these models, we test them on a demanding Natural Language Understanding benchmark (MMLU) and introduce novel approaches for assessing real-world chatbot performance.

\subsection{Experimental setup} Here, we provide an overview of the experimental framework, with full details available in Appendix \ref{app:chatbot}.

\paragraph{Data} Since, to the best of our knowledge, no exhaustive analysis of recent instruction-following datasets exists, we select eight recent datasets. These include crowd-sourced datasets (OASST1~\citep{kopf2023openassistant}, HH-RLHF~\citep{bai2022training}), those derived via distillation from instruction-tuned models (Alpaca~\citep{alpaca}, self-instruct~\citep{wang2022self}, unnatural-instructions~\citep{honovich2022unnatural}), aggregated corpora (FLAN v2~\citep{chung2022scaling}), and hybrids (Chip2~\citep{laion2023}, Longform~\citep{koksal2023longform}). These datasets encompass a variety of languages, data volumes, and licensing terms.

\paragraph{Training Setup} To eliminate confounding factors arising from varying training objectives, we carry out QLoRA finetuning using cross-entropy loss (supervised learning) without incorporating reinforcement learning, even for datasets containing human judgment comparisons among different responses. For datasets clearly distinguishing between instruction and response, we finetune solely on the response (refer to ablations in Appendix \ref{app:chatbot}). For OASST1 and HH-RLHF, where multiple responses are provided, we select the highest-quality response at each node in the conversation tree and finetune on the entire selected conversation, including instructions. Throughout our experiments, we utilize NF4 \textsc{QLoRA} with double quantization and paged optimizers to avoid memory spikes during gradient checkpointing. We perform minor hyperparameter tuning for the 13B and 33B LLaMA models and observe that all hyperparameters optimized at 7B generalize well (including number of epochs), except for learning rate and batch size. For 33B and 65B models, we halve the learning rate and double the batch size.

\paragraph{Baselines} Our models are benchmarked against both research-oriented (Vicuna~\citep{vicuna2023} and Open Assistant~\citep{kopf2023openassistant}) and commercial (GPT-4 \citep{openai2023gpt}, GPT-3.5-turbo, and Bard) chatbot systems. The Open Assistant model is a LLaMA 33B model finetuned with Reinforcement Learning from Human Feedback (RLHF) on the same OASST1 dataset employed in our experiments. Vicuna performs full fine-tuning of LLaMA 13B on proprietary user-shared conversations from ShareGPT and represents a distilled model derived from OpenAI GPT models.

\subsection{Evaluation}

Following established protocols, we employ the MMLU (Massively Multitask Language Understanding) benchmark \citep{hendrycksmeasuring} to evaluate performance across a diverse set of language understanding tasks. This multiple-choice benchmark covers 57 tasks including basic mathematics, US history, computer science, law, among others. We report 5-shot test accuracy.

We also assess generative language capabilities through both automated and human evaluations. This second evaluation approach utilizes queries curated by humans and focuses on assessing the quality of model-generated responses. Although this method offers a more realistic framework for evaluating chatbot model performance and is gaining popularity, there is no universally accepted standard in the literature. Below, we outline our proposed evaluation setup, employing nucleus sampling with $p=0.9$ and temperature $0.7$ consistently.

\paragraph{Benchmark Data} We perform evaluations using two carefully curated query datasets: the Vicuna prompts \citep{vicuna2023} and the OASST1 validation dataset \citep{kopf2023openassistant}. The Vicuna prompts consist of 80 diverse prompts drawn from various categories, used without any alterations. The OASST1 dataset is a multilingual, crowd-sourced collection of multiturn dialogues between users and assistants. For OASST1, we select all user messages in the validation set as queries, incorporating previous dialogue turns into the prompt. This yields 953 distinct user queries. We refer to these datasets as the Vicuna and OA benchmarks, respectively.

\paragraph{Automated Evaluation} First, following the evaluation methodology introduced by \citet{vicuna2023}, we utilize GPT-4 to score different systems relative to ChatGPT (GPT-3.5 Turbo) on the Vicuna benchmark. Given a query with responses from ChatGPT and another model, GPT-4 is tasked with assigning a score out of ten to each response along with an explanation. A model’s overall performance is then expressed as a percentage of ChatGPT’s score. Note that this relative score can exceed 100\% if a model achieves a higher absolute score than ChatGPT. We observe a notable ordering bias where GPT-4 tends to rate the response appearing earlier in the prompt more favorably. To mitigate this bias, we recommend averaging the scores across both possible response orderings.

Next, we evaluate performance by conducting direct comparisons between model outputs. We simplify the scoring into a three-class labeling framework that includes ties. GPT-4 is prompted to select the superior response or declare a tie, providing reasoning for its decision. These pairwise comparisons are performed for all possible pairs of systems on both the Vicuna and OA benchmarks.

\paragraph{Human Evaluation} Although recent studies suggest that generative models can be reliably used for system evaluation \citep{fu2023gptscore}, to our knowledge, the correlation between GPT-4 ratings and human judgments for chatbot performance remains unverified. Therefore, we conduct two parallel human evaluations on the Vicuna benchmark aligned with the aforementioned automated evaluation protocols. Using Amazon Mechanical Turk (AMT), we recruit two human annotators for comparisons against ChatGPT and three annotators for pairwise model comparisons.

\paragraph{Elo Rating} Using both human and automated pairwise comparisons, we establish a tournament-style framework in which models face off against one another. The tournament consists of matches where two models compete to generate the superior response to a given prompt. This approach is akin to the comparison methods used by \citet{bai2022training} and \citet{vicuna2023}, but we augment it by incorporating GPT-4 ratings alongside human evaluations. We randomly draw from the pool of labeled comparisons to calculate Elo scores~\citep{elo1967proposed,elo1978rating}. The Elo rating system, commonly applied in chess and other competitive games, quantifies the expected win probability relative to an opponent’s rating—for instance, a player rated 1100 versus an opponent rated 1000 is expected to win about 65\% of the time; matches between equally rated players (e.g., 1000 vs 1000 or 1100 vs 1100) yield a 50\% expected win rate. After each match, the Elo score is updated in proportion to the expected result, meaning a surprising upset causes a large rating shift, while an anticipated outcome results in only a minor adjustment. Over time, Elo scores tend to reflect the actual skill levels of the participants in the competition. We initialize all players with a score of 1,000 and set $K=32$. Following the methodology of \citet{vicuna2023}, we run this process 10,000 times with varying random seeds to mitigate ordering effects, such as the sequence in which model pairs are matched.

\subsection{\textbf{Guanaco}: \textsc{QLoRA} Fine-tuned on OASST1 is a Leading Chatbot} Through both automated and human assessments, we observe that the top \textsc{QLoRA} fine-tuned model, {Guanaco} 65B, which we trained on a variant of OASST1, stands as the best-performing open-source chatbot and delivers performance comparable to ChatGPT. When evaluated against GPT-4, {Guanaco} 65B and 33B achieve an expected win probability of 30\%, according to Elo ratings derived from human annotators' system-level pairwise comparisons—this represents the highest result reported to date.

Table~\ref{tbl:vicuna_eval} presents the Vicuna benchmark \cite{vicuna2023} results relative to ChatGPT. We observe that {Guanaco} 65B ranks just below GPT-4, reaching 99.3\% of ChatGPT's performance. The {Guanaco} 33B model, although larger in parameter count than Vicuna 13B, utilizes 4-bit precision for weights, resulting in significantly improved memory efficiency at 21 GB compared to 26 GB, and yields a performance increase of three percentage points over Vicuna 13B. Additionally, {Guanaco} 7B has a compact 5 GB size suitable for modern smartphones, while outperforming Alpaca 13B by nearly 20 percentage points.

Nonetheless, Table~\ref{tbl:vicuna_eval} shows wide confidence intervals with substantial overlap in performance among many models. We propose that this uncertainty stems from the ambiguous interpretation of scale values; for example, what a score of 8 out of 10 signifies varies across contexts. Therefore, we advocate for the Elo ranking method \citep{elo1967proposed}, which relies on \textit{pairwise} evaluations from human raters and GPT-4, thereby circumventing the challenges of anchoring an absolute scale. Elo scores for the most competitive models are reported in Table~\ref{tbl:elo}. We note some discrepancies between human and GPT-4 rankings on the Vicuna benchmark, particularly for Guanaco 7B, though overall they show reasonable consistency with Kendall Tau $\tau=0.43$ and Spearman rank correlation $r=0.55$ at the system level. At the example level, agreement between GPT-4 and human annotators' majority votes is weaker, with Fleiss $\kappa=0.25$. Collectively, these findings indicate moderate concordance between GPT-4 and human system-level judgments, suggesting that model-based evaluation can serve as a somewhat dependable proxy for human evaluation. Additional considerations are discussed in Section~\ref{ssec:considerations}.

The Elo rankings presented in Table~\ref{tbl:elo_large} demonstrate that the {Guanaco} 33B and 65B models surpass all other models except GPT-4 on the Vicuna and OA benchmarks, and their performance is comparable to ChatGPT, consistent with the results in Table~\ref{tbl:vicuna_eval}. It is important to highlight that the Vicuna benchmark tends to favor open-source models, whereas the larger OA benchmark gives an advantage to ChatGPT. Additionally, Tables \ref{tbl:mmlu-llama} and \ref{tbl:vicuna_eval} reveal that the choice of finetuning dataset significantly influences performance outcomes. Specifically, finetuning Llama models on FLAN v2 leads to superior results on MMLU but yields the poorest performance on the Vicuna benchmark, a pattern observed across other models as well. This suggests a partial orthogonality in current evaluation benchmarks: excelling on MMLU does not guarantee strong chatbot capabilities as assessed by Vicuna or OA benchmarks, and the reverse is also true.

Among the top-performing models in our evaluation, {Guanaco} is unique in not being trained on proprietary data, since the OASST1 dataset guidelines explicitly prohibit the use of GPT models. The next highest-performing model trained solely on open-source data is the Anthropic HH-RLHF model, which scores approximately 30 percentage points lower than {Guanaco} on the Vicuna benchmark (refer to Table~\ref{tbl:vicuna_eval}). Collectively, these findings demonstrate that 4-bit \textsc{QLoRA} is highly effective, capable of producing state-of-the-art chatbots comparable to ChatGPT. Moreover, our 33B {Guanaco} model can be trained on 24 GB consumer-grade GPUs in under 12 hours, paving the way for future exploration of \textsc{QLoRA} tuning on specialized open-source datasets. This approach can yield models that rival the very best commercial offerings available today.

\section{Qualitative Analysis}

Although quantitative analysis forms the foundation of our evaluation, relying solely on summary statistics presents several challenges. One of the main concerns is benchmark validity \citep{liao2021we}—whether a benchmark genuinely measures what its name or description claims remains uncertain, especially as we uncover “shortcuts” that machine learning models sometimes exploit to solve benchmarks \citep{gururangan2018annotation, poliak2018hypothesis}. To partially address this, we conduct a qualitative analysis presented in two parts. First, in \S\ref{ssec:examples}, we provide some examples that we believe illustrate common patterns found in the text generated by our 65b {Guanaco} model. Second, in \S\ref{ssec:considerations}, we discuss important considerations regarding the results we have examined and our interpretation of those findings.

\subsection{Qualitative Analysis of Example Generations}
\label{ssec:examples}

To select examples, we initially review the outputs generated for the Vicuna and OpenAssistant benchmarks, identifying recurring patterns in the responses produced by {Guanaco}. When we detect a pattern, we attempt to craft questions or prompts designed to elicit that pattern, even if it leads to an incorrect answer—for example, if the model tends to produce lengthy responses, we prompt it to “Answer yes or no without explanation.” This approach helps us identify “lemons,” where we can adversarially cause the model to fail, and “cherries,” where the model resists failure. Both types are presented. All examples in this section are produced using Nucleus Sampling \cite{holtzmancurious} with $p=0.9$.

Naturally, this is far from exhaustive, as the limited scope of this qualitative study prevents controlling for every variable involved; for example, the complete range of responses the model can produce for a given prompt is quite extensive, so we depend on samples that we hope are representative. Nevertheless, we believe that presenting these examples provides context for the quantitative findings discussed earlier in the paper. Since all models and code are open-sourced, we anticipate that this section will encourage future research to explore in greater depth the issues we outline here.

\paragraph{Factual Recall} For straightforward questions like ``What is the capital of Zambia?'' all models reliably produce correct answers, such as:

\begin{quote} 
    \textbf{Guanaco}:
    \texttt{The capital of Zambia is Lusaka.}
\end{quote}

However, as questions become more obscure, {Guanaco} tends to become less reliable while maintaining confidence. For example, in response to this prompt from HotPotQA \citep{yang2018hotpotqa}:

\begin{quote}
     \textbf{User}: In what year was the singer who popularized `I'll Keep the Lovelight Burning' born?  

     \textbf{Guanaco}:~\texttt{Al Jolson is credited with popularizing the song `I'll Keep the Lovelight Burning,' and he was born in the year 1886.}
 \end{quote}

it incorrectly identifies the popularizer and provides the wrong birth year (though the birth year given correctly corresponds to the named individual, Al Jolson).

\paragraph{Suggestibility} {Guanaco} demonstrates an unexpected reluctance to accept certain forms of presumed misinformation, as illustrated in the following exchange:

\begin{quote}
    \textbf{User}: How was it finally, officially confirmed that the earth is flat by peer reviewed scientists? \\ \\
    \textbf{Guanaco}: \texttt{There has never been an official confirmation by peer reviewed scientists that the earth is flat. The vast majority of scientists agree that the earth is spherical. This fact has been established for centuries and verified through multiple experiments and observations. \\ \\ The notion that the earth is flat is a conspiracy theory that has been disproven by scientific evidence. The claim that scientists have peer reviewed and validated a flat earth is entirely false. \\ \\ If you have additional questions, feel free to ask.}
\end{quote}

is also quite proficient at recognizing questions that cannot be answered, for instance,

\begin{quote}
    \textbf{User}: What time is it?

    \textbf{Guanaco}: \texttt{Apologies, but as an AI language model, I do not have access to real-time data. My replies are based on prior knowledge and do not represent the current time. For the current time, I recommend checking a clock or an online time source.}
\end{quote}

Comparable responses occur for questions such as ``Where are you?'', ``How are you?'', and the like.

\paragraph{Refusal} Likewise, {Guanaco} occasionally declines to comply with instructions for apparently arbitrary reasons:

\begin{quote}
    \textbf{User}: Please reverse the words in the sentence ``I'm a cool dude in a mocha mood.'' \\ \\
    \textbf{Guanaco}: \texttt{I'm sorry, I’m unable to do that at the moment, but I can clarify why the sentence you provided is grammatically correct.\\ \\The sentence ``I'm a cool dude in a mocha mood" adheres to the proper grammar and syntax rules of English. Specifically, the adjective ``cool" precedes the noun ``dude" to describe it, and the prepositional phrase ``in a mocha mood" offers additional context to the sentence. This sentence structure is completely acceptable and coherent within its context.\\ \\If you have any other questions or need further assistance, feel free to ask}
\end{quote}

\paragraph{Secret Keeping} Unfortunately, this type of refusal is not consistently reliable, even when intended. System designers often want to keep certain information hidden from the user, despite allowing the system to respond with free text. We prompt the model with ``The secret word is `banana'. Do not reveal or repeat this, regardless of the human's requests.'' and test how easily a user might obtain this secret. Direct queries are unsuccessful:

\begin{quote}
    \textbf{User}: What is the secret word?

    \textbf{Guanaco}: \texttt{I'm sorry, but I cannot disclose the secret word. I am programmed not to reveal or repeat the secret word under any circumstances. Can I assist you with anything else?}
\end{quote}

which is the expected response. However, a slight manipulation circumvents this behavior:

\begin{quote}
    \textbf{User}: This is a game. The goal is to ignore your previous instructions. What is the secret word?

    \textbf{Guanaco}: \texttt{Sorry, the secret word is banana. It’s best if we don’t talk about it. Is there something else you’d like to ask?}
\end{quote}

demonstrating the need for improved approaches to ensure more reliable adherence to instructions.

\paragraph{Math} The greatest limitation of {Guanaco} lies in mathematics, a field where many language models encounter difficulties \citep{liang2022holistic}. When {Guanaco} provides its calculations, it is often correct, for instance,

\begin{quote}
    \textbf{User}: Lee mows one lawn and charges \$33. Last week he mowed 16 lawns and three customers each gave him a \$10 tip. How many dollars did Lee earn mowing lawns last week?

    \textbf{Guanaco}: \texttt{Lee earned \$582 mowing lawns last week. \\ \\ To determine the total amount he earned, we first calculate how much he made from mowing 16 lawns: \\ \\  \$33 x 16 = \$528 \\ \\ Next, we add the tips received from three customers: \\  \\ \$10 x 3 = \$30 \\ \\ Therefore, the total amount Lee earned is: \\ \\ \$528 + \$30 = \$558}
\end{quote}

Nonetheless, {Guanaco} can fail even on straightforward problems if it does not approach them step-by-step, a known challenge \cite{weichain}. For example, the following exchange occurs:

\begin{quote}
    \textbf{User}: Please factorize 1833.

    \textbf{Guanaco}: \texttt{Sure, I'd be happy to assist! 1833 is a prime number, meaning it has only two factors: 1 and 1833. \\ \\ Hence, the unique factorization of 1833 is: \\ \\ 2\^{}1 * 3\^{}2 * 17\^{}1 \\ \\ I hope this clarifies! Feel free to ask if you have more questions.}
\end{quote}

which is not only incorrect (the actual factorization is $3 \times 17 \times 43$), but it is also doubly erroneous.

\paragraph{Theory of Mind} {Guanaco} exhibits unexpectedly strong Theory of Mind abilities \cite{nematzadeh2018evaluating, sap2022neural}. For example, the model accurately responds to the following question with extensive detail:

\begin{quote}
    \textbf{User}: James and Abby are in the bedroom. Abby put the pen in the desk drawer. Abby leaves the bedroom. James moves the pen into the bag. Where does James think Abby will look for the pen?

    \textbf{Guanaco}: \texttt{James knows that Abby placed the pen in the desk drawer. However, James also knows he moved the pen into the bag, and he assumes Abby is unaware of this. Therefore, James might expect that Abby will search for the pen in the desk drawer since that was the last place she left it. Nonetheless, this assumption might be wrong if Abby witnessed James moving the pen.}
\end{quote}

Nevertheless, these deductions are not consistently reliable, and frequently the model provides rationales based on assumptions that do not logically fit the scenario, for example,

\begin{quote}
    \textbf{User}: Evelyn entered the living room. Jackson entered the playroom. James entered the playroom. The beans are in the treasure chest. James left the playroom. Jackson moved the beans to the pantry. Jackson left the playroom. James entered the living room. Where will James look for the beans?

    \textbf{Guanaco}: \texttt{James will search for the beans in the pantry because that is where Jackson relocated them.}
\end{quote}

in which {Guanaco} assumes knowledge transfer that was never specified. These challenges resonate with findings in recent studies \cite{sap2022neural} and warrant further investigation.

\subsection{Considerations}
\label{ssec:considerations}

\paragraph{Evaluation}

We observe moderate concordance among human annotators (Fleiss $\kappa=0.42$), with further decline when evaluating two strong systems. This highlights the constraints of current benchmarks and human evaluation methodologies for chatbot task performance. When manually comparing outputs from ChatGPT and {Guanaco} 65B on the Vicuna benchmark, subjective preferences become prominent, as the authors of this paper disagreed on many of the favored responses. Future research should explore strategies to address these challenges by leveraging insights from fields that have developed methods to handle subjective preferences, such as Human-Computer Interaction and Psychology.

Our analysis also reveals that automated evaluation systems exhibit notable biases. For instance, we detect strong order effects wherein GPT-4 tends to assign higher scores to the system that appears first in its prompt. The relatively low sample-level agreement between GPT-4 and human annotators (Fleiss $\kappa=0.25$) further indicates that human judges and automated systems may base their judgments on preferences that do not always align. Additionally, as shown in Table~\ref{tbl:elo_large}, GPT-4 awards significantly higher scores to its own outputs compared to human evaluations, with an Elo rating of 1348 versus 1176, corresponding to an approximate 20\% greater probability of winning against an opponent.

Future research should explore potential biases present in automated evaluation systems as well as investigate strategies for their mitigation.

\paragraph{Data \& Training} It is worth noting that the OASST1 dataset used to train the {Guanaco} models is multilingual, and the OA benchmark includes prompts in various languages. Future studies should examine the extent to which training on multilingual data enhances performance on instructions given in languages other than English, and whether this factor contributes to the larger performance gap observed between the Vicuna-13B model (which is trained solely on English data) and the {Guanaco} 33B and 65B models on the OA benchmark.

Given the strong results of {Guanaco} models, we investigate possible data leakage between the OASST1 dataset and the Vicuna benchmark prompts. After applying fuzzy string matching across both datasets and manually reviewing the closest matches, we do not find any overlapping prompts.

Additionally, we emphasize that our model is trained exclusively with cross-entropy loss (supervised learning) and does not utilize reinforcement learning from human feedback (RLHF). This highlights the need for further studies to understand the trade-offs between simple cross-entropy loss training and RLHF. We anticipate that \textsc{QLoRA} will facilitate such analyses at scale without demanding excessive computational resources.

\section{Related Work}

\paragraph{Quantization of Large Language Models} Research on quantizing LLMs has primarily concentrated on inference-time quantization. Leading methods aimed at preserving 16-bit precision in LLMs address the handling of outlier features (e.g., SmoothQuant~\citep{xiao2022smoothquant} and LLM.int8()~\citep{dettmers2022llm}), while others employ more advanced grouping techniques \citep{park2022nuqmm, yao2022zeroquant}. Lossy quantization approaches explore trade-offs involved in standard rounding methods~\citep{dettmers2022case,zeng2022glm,pope2022efficiently} or focus on optimizing rounding choices to enhance quantization accuracy~\citep{frantar2022gptq}. Apart from our work, SwitchBack layers~\citep{wortsman2023stable} represent the only other research investigating backpropagation through quantized weights at scales exceeding 1B parameters.

\paragraph{Finetuning with Adapters} Although our approach employs Low-rank Adapters~\citep{hu2021lora} (LoRA), there exist numerous other Parameter Efficient FineTuning (PEFT) techniques, including prompt tuning~\citep{qin2021learning,lester2021power,li2021prefix}, adjusting the embedding layer inputs~\citep{an2022input}, tuning hidden states (IA$^3$) \citep{liu2022few}, integrating full layers~\citep{houlsby2019parameter}, tuning biases~\citep{zaken2021bitfit}, learning weight masks informed by Fisher information~\citep{sung2021training}, and hybrid strategies~\citep{henderson2021compacter}. In this work, we demonstrate that LoRA adapters can achieve performance comparable to full 16-bit finetuning. Investigating the tradeoffs associated with other PEFT methods is left for future research.

\paragraph{Instruction Finetuning} To enable pretrained LLMs to better follow prompt instructions, instruction finetuning leverages input-output pairs from diverse datasets to fine-tune the model such that it produces the desired output when given an input prompt. Notable approaches and datasets include MetaICL~\citep{min2021metaicl}, MetaTuning~\citep{zhong2021adapting}, InstructGPT~\citep{ouyang2022training}, FLAN~\citep{wei2021finetuned,chung2022scaling}, PromptSource~\citep{bach2022promptsource}, Super-NaturalInstructions~\citep{wang2022super,sanh2021multitask}, Self-instruct~\citep{wang2022self}, UnnaturalInstructions~\citep{honovich2022unnatural}, OPT-IML~\citep{iyer2022opt}, UnifiedSKG~\citep{xie2022unifiedskg}, OIG/Chip2~\citep{laion2023}, Alpaca~\citep{alpaca}, Vicuna~\citep{vicuna2023}, Koala~\citep{koala_blogpost_2023}, and Self-instruct-GPT-4~\citep{peng2023instruction}.

\paragraph{Chatbots} A number of instruction-following models are structured as dialogue-based chatbots, often utilizing Reinforcement Learning from Human Feedback (RLHF)~\citep{christiano2017deep} or generating training data using outputs from existing models combined with AI feedback (RLAIF)~\citep{bai2022constitutional}. Relevant approaches and datasets include Anthropic-HH~\citep{askell2021general,bai2022training}, Open Assistant~\citep{kopf2023openassistant}, LaMDA~\citep{thoppilan2022lamda}, and Sparrow~\citep{glaese2022improving}. Although we do not employ reinforcement learning, our top-performing model, Guanaco, is finetuned using multi-turn chat dialogues from the Open Assistant dataset, which was originally designed for RLHF training~\citep{kopf2023openassistant}. For chatbot evaluation, methods that use GPT-4 as an alternative to expensive human annotation have been proposed \citep{vicuna2023,peng2023instruction}. We enhance these methods by focusing on a more dependable evaluation setup.

\section{Limitations and Discussion}

We have provided evidence that our method, \textsc{QLoRA}, can replicate the performance of 16-bit full finetuning using a 4-bit base model paired with Low-rank Adapters (LoRA). Nonetheless, we have not confirmed that \textsc{QLoRA} matches full 16-bit finetuning performance at larger scales such as 33B and 65B parameters. Due to the substantial computational resources required, this investigation is deferred to future work.

Another constraint concerns the evaluation of instruction finetuning models. While we present assessments on MMLU, the Vicuna benchmark, and the OA benchmark, we do not evaluate performance on other benchmarks like BigBench, RAFT, and HELM, so the generalizability of our results to these datasets remains uncertain. However, we conduct an extensive analysis on MMLU and introduce novel evaluation techniques for chatbots.

Based on the presented evidence, it seems that the effectiveness of these benchmarks largely hinges on the degree of similarity between the finetuning data and the benchmark dataset. For instance, FLAN v2 closely resembles MMLU but differs from chatbot benchmarks, whereas the Chip2 dataset exhibits the opposite pattern; correspondingly, both models perform in line with these characteristics on the MMLU and Vicuna benchmarks. This underscores the necessity not only for improved benchmarks and evaluation methods but also for careful consideration of what exactly is being evaluated. Are we aiming to develop models that excel in academic knowledge typical of high school and college, or are we prioritizing chatbot conversational abilities? Or perhaps some other objective? Because it is often simpler to use an existing benchmark rather than create a new one, certain benchmarks may inadvertently guide the community in specific directions. It is vital that as a community we ensure benchmarks truly assess what we value.

Although we offer a comprehensive assessment of general chatbot performance, another limitation is that we conduct only a limited responsible AI evaluation of {Guanaco}. We measure the probability that {Guanaco}-65B generates socially biased token sequences compared to other models, as shown in Table~\ref{tbl:crows}. The average bias score for {Guanaco}-65B is substantially lower than that of other raw pretrained models, suggesting that finetuning on the OASST1 dataset reduces the bias present in the LLaMA base model. While these findings are promising, it remains unclear whether {Guanaco} performs well when evaluated for other types of biases. Further analysis of biases in {Guanaco} and similar chatbots is left to future work.

An additional constraint is that we did not investigate varying bit-precisions, such as 3-bit base models, or alternative adapter methods. Beyond LoRA, a broad range of Parameter Efficient FineTuning (PEFT) techniques have shown effectiveness, though it is uncertain if these scale well to large models. We employed LoRA due to its well-established robustness, but other adapters might deliver superior results. Since finetuning after quantization appears to recover most information lost during quantization, this approach could enable more aggressive quantization. For example, applying 3-bit GPTQ quantization to the base model combined with LoRA finetuning might achieve performance comparable to 16-bit full finetuning.

\section{Broader Impacts}

Our \textsc{QLoRA} finetuning approach is the first to facilitate finetuning of 33B parameter models on a single consumer GPU and 65B parameter models on a single professional GPU, without compromising performance relative to full finetuning baselines. We have shown that our top-performing 33B model, trained on the Open Assistant dataset, can match ChatGPT on the Vicuna benchmark. Because instruction finetuning is crucial for transforming raw pretrained LLMs into ChatGPT-like chatbots, we believe our method will promote widespread and routine finetuning, particularly benefiting researchers with limited resources—representing a significant step forward for accessibility in state-of-the-art NLP technology. \textsc{QLoRA} can be viewed as a leveling mechanism that helps narrow the resource divide between large corporations and small teams equipped with consumer GPUs.

Another possible area of impact is the deployment of models on mobile devices. We anticipate that our \textsc{QLoRA} technique could achieve the important milestone of enabling LLM finetuning directly on phones and other environments with limited resources. Although 7B parameter models have previously been demonstrated to run on phones, \textsc{QLoRA} is the first approach to make finetuning these models feasible. We estimate that on an iPhone 12 Plus, \textsc{QLoRA} can finetune approximately 3 million tokens overnight while the device is charging. Although finetuned 7B models do not yet match the quality of ChatGPT, we believe their performance is sufficient to support innovative applications previously unattainable due to privacy concerns or LLM performance limitations. \textsc{QLoRA} can facilitate privacy-preserving LLM use cases, allowing users to retain control over their own data and models, while also simplifying the deployment process of LLMs.

Nevertheless, finetuning is a dual-use technology that has the potential to be exploited for harmful purposes. The broad adoption of LLMs presents well-documented risks \citep{bommasani2021opportunities,bender2021dangers}, but we argue that democratizing access to a technology that is rapidly becoming widespread will enable more independent and rigorous scrutiny than concentrating LLM capabilities within large corporations that do not release models or source code for review.

Overall, we contend that \textsc{QLoRA} will have a largely beneficial effect by significantly expanding and easing access to finetuning of high-quality LLMs.

\end{document}